\section{Dream Record / 梦境记录}

Hmm. I had a dream, and dreams, of course, are astonishingly free.

\begin{translationbox}
Hmm，我做了一个梦，而梦境这种东西，本来就自由得惊人。
\end{translationbox}

The more interesting first half has already dissolved into mist. What remains is the ending: I was running toward a gate formed by an enormous split in a campus fence, and by slipping through it I escaped one world and entered another.

\begin{translationbox}
梦境更有意思的前半段已经模糊成雾了。真正留下来的，是最后的那一幕：我一路狂奔，冲向校园围栏上一道巨大裂缝形成的门；穿过去的瞬间，我从原本的世界逃离，跌进了另一个世界。
\end{translationbox}

\begin{figure}[H]
  \centering
  \begin{subfigure}[t]{0.47\textwidth}
    \centering
    \includegraphics[width=\linewidth]{1775240989443.png}
    \caption{逃离前的压迫感场景}
  \end{subfigure}
  \hfill
  \begin{subfigure}[t]{0.47\textwidth}
    \centering
    \includegraphics[width=\linewidth]{1775241009838.png}
    \caption{穿过门缝后的过渡瞬间}
  \end{subfigure}
  \caption{梦境前半段残留的视觉切片}
\end{figure}

In that world, every resource had to be obtained by entering a world behind a door, and each person possessed a door-world of their own.

\begin{translationbox}
在那个世界里，一切资源都必须通过“进入门后的世界”才能获得，而且每个人都有一个只属于自己的门世界。
\end{translationbox}

The one I entered was eerie in a particularly intimate way. It looked like my hometown, yet everything had been enlarged beyond proportion. The buildings were vastly bigger than they should have been, so familiarity and estrangement pressed against each other at the same time.

\begin{translationbox}
我进入的那个门世界，有一种非常私人的诡异感。它看起来像老家，却又被夸张地放大了；建筑比现实中巨大得多，于是熟悉感和异样感同时压了过来。
\end{translationbox}

\begin{scenequote}
The door-world was not an unknown fantasy realm. It was a private memory-space, stretched, distorted, and lightly gamified.
That half-familiar structure made the fear sharper because it never felt entirely foreign.
\end{scenequote}

\begin{figure}[H]
  \centering
  \begin{subfigure}[t]{0.47\textwidth}
    \centering
    \includegraphics[width=\linewidth]{1775241628999.png}
    \caption{巨大化的老家式建筑}
  \end{subfigure}
  \hfill
  \begin{subfigure}[t]{0.47\textwidth}
    \centering
    \includegraphics[width=\linewidth]{1775241783883.png}
    \caption{古旧庭院与危险逼近}
  \end{subfigure}
  \caption{门世界的核心空间}
\end{figure}

I entered a household that resembled an old family courtyard. I can no longer remember the monsters clearly, but I remember the method of resistance with absurd precision: I swung a pole used to support bedding in the sun and used it to kill them one after another.

\begin{translationbox}
我走进一户人家，像是那种老式家宅风格的庭院。怪物的具体样子我已经记不清了，但我却异常清楚地记得自己是怎么反击的：我抓起一个晾晒被褥用的栏杆支架，挥舞着把那些怪物一个个击杀。
\end{translationbox}

As the bodies accumulated, it seemed I had attracted the notice of something far stronger. What fixed its attention on me may have been a dragon, or some other immense creature. In any case, it unleashed a force so powerful that I was lifted into the air inside a vortex, suspended high above the ground.

\begin{translationbox}
随着倒下的怪物越来越多，我似乎也惊动了一个更强大的存在。盯上我的东西也许是一条龙，也许是别的什么庞然巨物。总之，它释放出一种极强的吸力，把我卷进一股龙卷风里，整个人被拽上高空，悬在半空中。
\end{translationbox}

My companion covered the base of the whirlwind with something from the ground, and the vortex collapsed. Then we ran. We fled through one street and then another, while the scenery kept mutating around us, until at last we found ourselves near the same gate through which we had first entered the world.

\begin{translationbox}
伙伴在地面上用某种东西盖住了龙卷风的根部，于是那股旋风一下子散掉了。接着我们开始逃命，从一条街跑到另一条街，四周场景不断变换，直到最后终于逼近我们最初进入这个世界时的那道门。
\end{translationbox}

But when we reached the exit, I was stuck. My companion pulled at me, yet I would not come free. On the gate there was a condition for departure, something like an inequality that had to be satisfied before the world would release me.

\begin{translationbox}
可真正到了出口，我却被卡住了。伙伴拼命拉我，我还是出不去。门上写着离开这个世界的条件，像是一道只有“任务完成后”才能满足的不等式判定。
\end{translationbox}

In the end we resorted to something like a reward hack. I changed the name to `xx < yy`, and when the world reached its settlement moment, the system misread me as having fulfilled the condition. That was how I escaped.

\begin{translationbox}
最后我们用了一个近乎奖励漏洞的办法。我把名称改成了 `xx < yy`，于是当门世界的结算时刻到来时，系统把我误判成已经满足条件，我就这样成功逃了出去。
\end{translationbox}

\section{Why Dreams Feel More Real Than Games / 为什么梦比游戏更真}

\begin{reflectionbox}
Dreams are more overwhelming than most virtual entertainments not merely because they are vivid, but because they seal off memory, mobilize emotion, and alter the environment in response to us.
Inside a dream, we do not know that it is false, so every fear, sprint, decision, and act of resistance comes from a less-defended self.
\end{reflectionbox}

I think of visiting haunted attractions at a theme park with younger relatives. We all know the performance is fake, and because we know it, our emotions never fully surrender. The same limit appears in many games. Whether it is an infection mode, a challenge mode, or an asymmetrical horror title, we are still ultimately facing a flat screen. The immersion can be real, but it is never complete.

\begin{translationbox}
我会想起以前和侄子、外甥去欢乐谷逛恐怖乐园。因为大家都知道那是假的，所以情绪始终很难真正交出去。很多游戏也是一样。无论是某种生化模式、挑战模式，还是非对称恐怖游戏，说到底我们面对的仍是一块屏幕。沉浸感当然会有，但那种沉浸总归不彻底。
\end{translationbox}

Dreams are different. In a dream, the world temporarily strips away its own label of unreality, and everything presents itself as immediate fact. Motion, combat, sensation, character, danger: all of it is intensified into a full-bodied experience. More importantly, the dream keeps reshaping itself in response to what we do, so our reactions do not merely happen inside the world; they help generate it.

\begin{translationbox}
梦不一样。梦里的世界会暂时抹掉自己“并不真实”的标签，于是一切都以“正在发生”的方式压到我们身上。动作、搏斗、感官、人物、危险，都会被放大成一种全身性的即时体验。更重要的是，梦境还会随着我们的反应不断改写自己，所以我们不只是身处其中，也在参与生成它。
\end{translationbox}

For that reason, dreams feel like an intensely private immersive medium. They project the inner world, give shape to the subconscious, and release compressed emotion all at once. Situations that would be impossible to inhabit so fully in waking life can, in dreams, be lived almost sensorily from the inside.

\begin{translationbox}
正因为如此，梦更像一种高度私人化的沉浸式媒介。它既是内心世界的投射，也是潜意识的具象化，同时还是被压缩情绪的一次集中释放。那些在清醒生活中几乎不可能被完整经历的处境，在梦里却能以接近全感官的方式，从内部被真正活过一次。
\end{translationbox}

\section{Condensed Insights / 核心观点}

\begin{itemize}[leftmargin=1.5em]
  \item A dream feels real less because it is vivid than because, inside it, we forget that it is virtual.
  \item The door-world suggests that the subconscious can rebuild private memory into a traversable, rule-bound system.
  \item Reactions inside dreams may reveal instinct more directly than many declarations made while awake.
  \item Compared with conventional entertainment, dreams feel deeper because they behave like environments that answer back.
\end{itemize}

\begin{translationbox}
\begin{itemize}[leftmargin=1.5em]
  \item 梦之所以显得真实，不只是因为它够逼真，更因为人在梦里会暂时忘记它的虚拟性。
  \item “门世界”说明潜意识会把私人记忆重新组装成一个可探索、可闯关、也可逃离的系统。
  \item 梦中的反应，往往比清醒时经过修饰的表达更接近本能。
  \item 相比一般娱乐媒介，梦更像一个会即时回应主体行为的动态环境，因此沉浸得更深。
\end{itemize}
\end{translationbox}

\section{Virtual Comment Thread / 虚拟评论区}

The following three comments simulate close-friend responses from upper-middle-class, college-educated readers aged roughly eighteen to twenty-six.
They are written in informal English first, then rendered into polished Chinese, so the page reads like an actual social post with a bilingual comment thread.

\begin{translationbox}
下面三则短评模拟的是 18--26 岁之间、中上层背景、至少受过本科教育的熟人留言。
呈现方式仍然是“英文原评论在前，中文润色译文在后”，让这一页更像真实社交媒体里的双语评论区。
\end{translationbox}

\begin{commentbox}{评论一}
Bro wait, this is lowkey the coolest dream lore dump ever T\_T the campus gate turning into your own door-world and then you literally code-hacked the ending? Nah, dreams really do clear games sometimes because your brain gives it full surround-sound reality. Pls write the sequel when you sleep again.
\end{commentbox}

\begin{translationbox}
Bro，等等，这简直是我最近看到最酷的梦境设定帖了 T\_T 校园围栏一下子变成你自己的门世界，最后你居然还是靠“代码式作弊”把结局改出来的？说真的，有时候梦比游戏还强，就是因为大脑会把那种全景式的真实感直接灌满。拜托你下次做梦记得继续更后续。
\end{translationbox}

\begin{commentbox}{评论二}
ok this is insanely cinematic?? the campus gate turning into your own private door-world, then suddenly monster-fighting in giant old-home physics, then literally hacking the rules to escape... that dream brain of yours is unmatched ;\_; dreams really do hit different because they feel 100\% true while you're inside them
\end{commentbox}

\begin{translationbox}
ok，这个真的太有电影感了吧？校园那道门一下子转成你自己的私人门世界，然后又突然切到“巨大老家物理系统”里的打怪逃生，最后还直接黑进规则里跑路……你这个做梦能力真没人比得上 ;\_; 梦之所以和别的不一样，就是因为人在里面的时候，会百分之百相信那一切真的正在发生。
\end{translationbox}

\begin{commentbox}{评论三}
This is so violently you though :') escaping through a campus gate into your own private door-world, then ending up monster-fighting in giant old-home mode, then literally math-hacking your way out?? No bc your brain makes better immersive sims than half the games you send me lol T\_T
\end{commentbox}

\begin{translationbox}
但这也真的太像你了吧 :') 先是从校园围栏逃进自己的私人门世界，然后又掉进“巨大化老家地图”里开始打怪，最后甚至还靠数学式 hack 把自己搞出去？？不夸张地说，你这脑子造出来的沉浸模拟，已经比你平时发给我的一半游戏都强了 lol T\_T
\end{translationbox}
